\documentclass[10pt,a4paper]{amsart}
\usepackage[margin=19mm]{geometry}
\usepackage{amsmath,amssymb,mathtools}
\usepackage[scr=rsfs,cal=euler]{mathalfa}
\usepackage{xfrac}
\usepackage[unicode,hidelinks,bookmarks=false]{hyperref}
\newcommand{\ie}{i.e.\ }
\newcommand{\cf}{cf.\ }
\newcommand{\resp}{respectively\ }
\newcommand{\Subsection}{\subsection}
\providecommand{\nobreakdash}{\nobreak\hbox{-}}
\setlength{\emergencystretch}{2em}
\multlinegap0pt
\usepackage{bm}
\usepackage{bbm}
\usepackage{dsfont}
\usepackage{xspace}
\usepackage{cancel}
\usepackage{mathtools}
\usepackage{amsthm}
\makeatletter
\@ifundefined{theorem}{\newtheorem{theorem}{Theorem}}{}
\@ifundefined{lemma}{\newtheorem{lemma}[theorem]{Lemma}}{}
\@ifundefined{proposition}{\newtheorem{proposition}[theorem]{Proposition}}{}
\makeatother
\makeatletter
\expandafter\ifx\csname tcb\endcsname\relax
\newenvironment{tcb}
{\color{blue}}
\fi
\makeatother
\title[Torsional Rigidity and Spherical Deficit for a Dirichlet Pro]{Torsional Rigidity and Spherical Deficit for a Dirichlet Problem on Riemannian Manifolds}
\author{Maria Andrade, Allan Freitas}
\date{Source: arXiv:2605.30088v1 (CC BY 4.0)}
\begin{document}
\maketitle

\noindent\emph{Source and licence.} Torsional Rigidity and Spherical Deficit for a Dirichlet Problem on Riemannian Manifolds. Version arXiv:2605.30088v1 (CC BY 4.0), distributed under the Creative Commons Attribution 4.0 International licence, as recorded in the arXiv licence metadata for this exact version and shown on the abstract page https://arxiv.org/abs/2605.30088v1. Full rights notice: licenses/arxiv-2605.30088-CC-BY-4.0.txt.\par\smallskip

\noindent\textit{Editorial scope.} Counted unit: the Proposition whose source label is \texttt{intlemma} in the article (source0.tex lines 668--688, source label \texttt{intlemma}), delivered together with the article's own complete original proof of it (source0.tex lines 690--1490, one \texttt{proof} environment of 801 lines). No mathematics is written, repaired or re-derived: statement and proof are the article's own text line for line. Required context retained uncounted: serrinoka (lines 241--244); lemma1 (lines 289--297); Lapl\_V (lines 291--295); deltaphi1 (lines 598--602). Every definition, display and label of those lines is retained unchanged. Omitted from this package and not redistributed: 1--240, 245--288, 296--597, 603--667, 689--689, 1491--1778. Nothing inside a retained excerpt is omitted or altered: every retained body is delivered exactly as the article's lines are written, so no omission receipt of any kind accompanies this package. Citations: the retained excerpts carry no citation command at all, so no bibliography entry of the article is transcribed and no reference is redistributed. Reference adaptation: none was needed and none was made; every reference inside the delivered unit resolves inside it, so no unresolved reference or citation is produced. Printed slips kept literal: no printed word, symbol, hypothesis or number of the counted statement or of its proof was altered. Layout and class: the article's own class is \texttt{amsart} with packages including hyperref, geometry, amsmath, amssymb, latexsym, soul, cite, mathrsfs, color, enumitem, graphicx, enumerate, cancel, babel; the wrapper is the lane's pinned \texttt{amsart} editorial wrapper at 10pt on A4, which restates the theorem-like environments the retained text needs and adds the article's own macro definitions that the retained text needs (none), with extra wrapper packages (bm, bbm, dsfont, xspace, cancel, mathtools). The delivered numbering is the unit's own. Assets: no retained span contains a figure environment, a graphics-inclusion command, a PDF-page inclusion, a table or a TikZ picture; no image file is copied into this package, the package declares no asset dependency and no image file is redistributed with it. Licence: version arXiv:2605.30088v1 of this article is distributed under the Creative Commons Attribution 4.0 International licence, as recorded in the arXiv licence metadata for this exact version and shown on the abstract page https://arxiv.org/abs/2605.30088v1. A full rights notice is delivered at licenses/arxiv-2605.30088-CC-BY-4.0.txt. Characters: every retained excerpt is pure ASCII, so the delivered engine, XeLaTeX with its default Latin Modern text font, reports no missing character.
\begin{eqnarray}
\Delta u + nku  &=& -1 \quad \text{in} \quad \Omega, \quad u = 0 \quad \text{on} \quad \partial \Omega, \label{serrinoka}\\
u_{\nu} &=& -c, \quad \text{on} \quad \partial \Omega, \label{serrinokb}    
\end{eqnarray}

\begin{lemma}\label{lemma1}
Let $(M^{n},g)$ be a Riemannian manifold endowed with a conformal vector field $X$ with conformal factor $V$. Then
\begin{equation}\label{Lapl_V}
-(n-1)\Delta V
=
\frac{1}{2}X(R)+VR,
\end{equation}
where $R$ denotes the scalar curvature of $(M^{n},g)$.
\end{lemma}

\begin{equation}\label{deltaphi1}
\nabla^{2}\varphi
=
-\left(\frac{1}{n}+K\varphi\right)g
\end{equation}

\begin{proposition}[General Integral Identity]\label{intlemma}
Let $(M^n, g)$ be a Riemannian manifold endowed with a conformal vector field $X$ with conformal factor $V$. Let $\Omega$ be a bounded domain with $C^2$ boundary, and let $u$ be a solution of
\[
\Delta u=-(1+nku), \qquad k\in\mathbb{R},
\]
in $\Omega$. Assume further that $\varphi$ is a solution of \eqref{deltaphi1}. Then the following identity holds:
\begin{eqnarray}\label{eqlema}
\displaystyle\int_{\Omega}-Vu\Delta Pdv&=&\displaystyle\int_{\partial \Omega}|\nabla u|^2[V(u-\varphi)_{\nu}-(u-\varphi)V_{\nu}]d\sigma
+2\displaystyle\int_{\partial \Omega}u\nabla^2u((u-\varphi)\nabla V-V\nabla (u-\varphi),\nu)d\sigma\nonumber\\
&&+k\displaystyle\int_{\partial \Omega}u^2[(u-\varphi)V_{\nu}-V(u-\varphi)_{\nu}]d\sigma
-\dfrac{2}{n(n-1)}\displaystyle\int_{\Omega}Vu(u-\varphi)(R-n(n-1)k)dv\nonumber\\
&&-2\displaystyle\int_{\Omega}u(u-\varphi)(\operatorname{Ric}-(n-1)kg)(\nabla u,\nabla V)dv
-2\displaystyle\int_{\Omega}Vu(\operatorname{Ric}-(n-1)kg)(\nabla u,\nabla \varphi)dv\nonumber\\
&&-\dfrac{k}{n-1}\displaystyle\int_{\Omega}Vu^2(u-\varphi)(R-(n-1)nk)dv
-\dfrac{1}{n-1}\displaystyle\int_{\Omega}V|\nabla u|^2(u-\varphi)(R-(n-1)nk)dv\nonumber\\
&&-\dfrac{k}{2(n-1)}\displaystyle\int_{\Omega}u^2(u-\varphi)X(R)dv
-\dfrac{1}{2(n-1)}\displaystyle\int_{\Omega}X(R)|\nabla u|^2(u-\varphi)dv\nonumber\\
&&-\dfrac{1}{n(n-1)}\displaystyle\int_{\Omega}u(u-\varphi)X(R)dv
-2\displaystyle\int_{\Omega}u(u-\varphi)\langle\nabla^2u,\mathring{\nabla}^{2}V\rangle dv.
\end{eqnarray}
\end{proposition}

\begin{proof} 
The proof is a long calculation.
Initially, observe that
\[
\operatorname{div}\bigl(Vu\nabla P-P\nabla(Vu)\bigr)
=
Vu\,\Delta P-P\,\Delta(Vu).
\]
Integrating this identity over \(\Omega\) and applying the divergence theorem, we obtain
\begin{equation}\label{eqin}
-\int_{\Omega}Vu\,\Delta Pdv
=
-\int_{\partial\Omega}Vu\,P_{\nu}dv
+\int_{\partial\Omega}P\,(Vu)_{\nu}dv
-\int_{\Omega}P\,\Delta(Vu)dv.
\end{equation}

Next, using the equation \(\Delta u=-(1+nku)\) together with Lemma \ref{lemma1}, we compute
\begin{align}\label{eqnablaVu}
\Delta(Vu)
&=
V\Delta u+u\Delta V+2\langle\nabla V,\nabla u\rangle \nonumber\\
&=
-V(1+nku)-\frac{uX(R)}{2(n-1)}-\frac{R}{n-1}Vu
+2\langle\nabla V,\nabla u\rangle \nonumber\\
&=
-V-\frac{Vu}{n-1}\bigl(R+n(n-1)k\bigr)
+2\langle\nabla V,\nabla u\rangle
-\frac{uX(R)}{2(n-1)}.
\end{align}

Substituting \eqref{eqnablaVu} into \eqref{eqin}, we arrive at
\begin{align}\label{eq43}
-\int_{\Omega}Vu\,\Delta Pdv
={}&
-\int_{\partial\Omega}Vu\,P_{\nu}d\sigma
+\int_{\partial\Omega}P\,(Vu)_{\nu}d\sigma
+\int_{\Omega}PVdv \nonumber\\
&\quad
+\frac{1}{n-1}\int_{\Omega}PVu\bigl(R+n(n-1)k\bigr)dv
-2\int_{\Omega}P\langle \nabla V,\nabla u\rangle dv\nonumber\\
&\quad
+\frac{1}{2(n-1)}\int_{\Omega}Pu\,X(R)dv.
\end{align}

We now make use of the auxiliary function \(\varphi\). Applying integration by parts, we obtain
\begin{align}\label{eq44}
\int_{\partial\Omega}P(Vu)_{\nu}d\sigma
&=
\int_{\partial\Omega}P\bigl(V(u-\varphi)\bigr)_{\nu}d\sigma
+
\int_{\partial\Omega}P(\varphi V)_{\nu}d\sigma
\nonumber\\
&=
\int_{\partial\Omega}PV(u_{\nu}-\varphi_{\nu})d\sigma
+
\int_{\partial\Omega}Pu\,V_{\nu}d\sigma
+
\int_{\partial\Omega}PV\varphi_{\nu}d\sigma
\nonumber\\
&=
\int_{\partial\Omega}PV(u_{\nu}-\varphi_{\nu})d\sigma
+
\int_{\Omega}\operatorname{div}(uP\nabla V)dv
+
\int_{\partial\Omega}PV\varphi_{\nu}d\sigma
\nonumber\\
&=
\int_{\partial\Omega}PV(u_{\nu}-\varphi_{\nu})d\sigma
+
\int_{\Omega}P\langle\nabla u,\nabla V\rangle dv
+
\int_{\Omega}u\langle\nabla P,\nabla V\rangle dv
\nonumber\\
&\quad
+
\int_{\Omega}uP\,\Delta Vdv
+
\int_{\partial\Omega}PV\varphi_{\nu}d\sigma
\nonumber\\
&=
\int_{\partial\Omega}PV(u_{\nu}-\varphi_{\nu})d\sigma
+
\int_{\Omega}P\langle\nabla u,\nabla V\rangle dv
+
\int_{\Omega}u\langle\nabla P,\nabla V\rangle dv
\nonumber\\
&\quad
-
\frac{1}{n-1}
\int_{\Omega}uP\left(\frac{X(R)}{2}+VR\right)dv
+
\int_{\partial\Omega}PV\varphi_{\nu}d\sigma.
\end{align}

Next, integrating by parts once again, we obtain
\begin{align}\label{eq45}
\int_{\partial\Omega}VP\varphi_{\nu}d\sigma
&=
\int_{\Omega}\operatorname{div}(VP\nabla\varphi)dv
\nonumber\\
&=
\int_{\Omega}VP\,\Delta\varphi dv
+
\int_{\Omega}P\langle\nabla V,\nabla\varphi\rangle dv
+
\int_{\Omega}V\langle\nabla P,\nabla\varphi\rangle dv
\nonumber\\
&=
-\int_{\Omega}VPdv
-nk\int_{\Omega}VP\varphi dv
+
\int_{\Omega}P\langle\nabla V,\nabla\varphi\rangle dv
+
\int_{\Omega}V\langle\nabla P,\nabla\varphi\rangle dv,
\end{align}
where we used the identity \(\Delta\varphi=-(1+nk\varphi)\). On the other hand, recalling that the \(P\)-function is given by
\begin{equation}\label{P_func1}
P=|\nabla u|^{2}+\frac{2}{n}u+ku^{2},    
\end{equation}
we compute
\begin{equation}\label{P_funct2}
\nabla P
=
2\nabla^{2}u(\nabla u,\cdot)
+\frac{2}{n}\nabla u
+2ku\nabla u.    
\end{equation}
Substituting these expressions into \eqref{eq45}, we arrive at
\begin{align}\label{eq46}
\int_{\partial\Omega}VP\varphi_{\nu} d\sigma
&=
-\int_{\Omega}VP dv
-nk\int_{\Omega}VP\varphi dv
+
\int_{\Omega}P\langle\nabla V,\nabla\varphi\rangle dv
\nonumber\\
&\quad
+
2\int_{\Omega}V\nabla^{2}u(\nabla u,\nabla\varphi)dv
+
\frac{2}{n}\int_{\Omega}V\langle\nabla u,\nabla\varphi\rangle dv
\nonumber\\
&\quad
+
2k\int_{\Omega}uV\langle\nabla u,\nabla\varphi\rangle dv.
\end{align}

In order to handle the Hessian term appearing in \eqref{eq46}, we recall the Ricci identity
\begin{equation}\label{Ricci_id}
\operatorname{div}(\nabla^2u)
=
\nabla(\Delta u)+\operatorname{Ric}(\nabla u,\cdot).    
\end{equation}
Using this formula, we compute
\begin{align}\label{eq_auxiliar}
\operatorname{div}\bigl(Vu\nabla_{\nabla\varphi}\nabla u\bigr)
&=
\operatorname{div}\bigl(Vu\,\nabla^{2}u(\nabla\varphi,\cdot)\bigr)
\nonumber\\
&=
u\nabla^{2}u(\nabla V,\nabla\varphi)
+
V\langle\nabla u,\nabla_{\nabla\varphi}\nabla u\rangle
+
Vu\,\operatorname{div}(\nabla_{\nabla\varphi}\nabla u)
\nonumber\\
&=
u\nabla^{2}u(\nabla V,\nabla\varphi)
+
V\nabla^{2}u(\nabla u,\nabla\varphi)
+
Vu\,\operatorname{Ric}(\nabla u,\nabla\varphi)
\nonumber\\
&\quad
+
Vu\langle\nabla\Delta u,\nabla\varphi\rangle
+
Vu\langle\nabla^{2}u,\nabla^{2}\varphi\rangle.
\end{align}

Integrating \eqref{eq_auxiliar} over \(\Omega\) and using the divergence theorem, we deduce
\begin{align}\label{eq47}
\int_{\Omega}V\nabla^{2}u(\nabla u,\nabla\varphi)dv
&=
\int_{\partial\Omega}Vu\,\nabla^{2}u(\nabla\varphi,\nu)d\sigma
-
\int_{\Omega}u\nabla^{2}u(\nabla V,\nabla\varphi)dv
\nonumber\\
&\quad
-
\int_{\Omega}Vu(\operatorname{Ric}-nkg)(\nabla u,\nabla\varphi)dv
-
\int_{\Omega}Vu\langle\nabla^{2}u,\nabla^{2}\varphi\rangle dv,
\end{align}
where we have used \(\nabla\Delta u=-nk\nabla u\).

Since \(\varphi\) and \(u\) satisfy \eqref{deltaphi1} and \eqref{serrinoka}, respectively, we infer that
\begin{align}
\int_{\Omega}Vu\langle\nabla^{2}u,\nabla^{2}\varphi\rangle dv
&=
\int_{\Omega}Vu\left[-\left(\frac{1}{n}+k\varphi\right)\Delta u\right]dv
\nonumber\\
&=
\frac{1}{n}\int_{\Omega}Vu\,\Delta\varphi\,(-1-nku)dv
\nonumber\\
&=
-\frac{1}{n}\int_{\Omega}Vu\,\Delta\varphi dv
-k\int_{\Omega}Vu^{2}\Delta\varphi dv.
\nonumber
\end{align}

Applying the divergence theorem once again, together with the identities
\[
\operatorname{div}(Vu\nabla\varphi)
=
u\langle\nabla V,\nabla\varphi\rangle
+
V\langle\nabla u,\nabla\varphi\rangle
+
Vu\Delta\varphi,
\]
and
\[
\operatorname{div}(Vu^{2}\nabla\varphi)
=
\langle\nabla(Vu^{2}),\nabla\varphi\rangle
+
Vu^{2}\Delta\varphi,
\]
we obtain
\begin{align}
\int_{\Omega}Vu\langle\nabla^{2}u,\nabla^{2}\varphi\rangle dv
&=
\frac{1}{n}\int_{\Omega}u\langle\nabla V,\nabla\varphi\rangle dv
+
\frac{1}{n}\int_{\Omega}V\langle\nabla u,\nabla\varphi\rangle dv
\nonumber\\
&\quad
-
\frac{1}{n}\int_{\partial\Omega}Vu\,\varphi_{\nu}d\sigma
+
k\int_{\Omega}\langle\nabla(Vu^{2}),\nabla\varphi\rangle dv
-
k\int_{\partial\Omega}Vu^{2}\varphi_{\nu}d\sigma.
\nonumber
\end{align}

Proceeding further, we arrive at
\begin{align}\label{eq48}
\int_{\Omega}Vu\langle\nabla^{2}u,\nabla^{2}\varphi\rangle dv
&=
-\frac{1}{n}\int_{\partial\Omega}Vu\,\varphi_{\nu}d\sigma
+
\frac{1}{n}\int_{\Omega}u\langle\nabla V,\nabla\varphi\rangle dv
+
\frac{1}{n}\int_{\Omega}V\langle\nabla u,\nabla\varphi\rangle dv
\nonumber\\
&\quad
+
2k\int_{\Omega}Vu\langle\nabla u,\nabla\varphi\rangle dv
+
k\int_{\Omega}u^{2}\langle\nabla V,\nabla\varphi\rangle dv
-
k\int_{\partial\Omega}Vu^{2}\varphi_{\nu}d\sigma.
\end{align}

Plugging this into \eqref{eq47}, we obtain
\begin{align}\label{eqstar}
\int_{\Omega}V\nabla^{2}u(\nabla u,\nabla\varphi)dv
&=
\int_{\partial\Omega}Vu\,\nabla^{2}u(\nabla\varphi,\nu)d\sigma
-
\int_{\Omega}u\nabla^{2}u(\nabla V,\nabla\varphi)dv
\nonumber\\
&\quad
-
\int_{\Omega}Vu(\operatorname{Ric}-nkg)(\nabla u,\nabla\varphi)dv
+
\frac{1}{n}\int_{\partial\Omega}Vu\,\varphi_{\nu}d\sigma
\nonumber\\
&\quad
-
\frac{1}{n}\int_{\Omega}u\langle\nabla V,\nabla\varphi\rangle dv
-
\frac{1}{n}\int_{\Omega}V\langle\nabla u,\nabla\varphi\rangle dv
\nonumber\\
&\quad
+
k\int_{\partial\Omega}Vu^{2}\varphi_{\nu}d\sigma
-
2k\int_{\Omega}Vu\langle\nabla u,\nabla\varphi\rangle dv
\nonumber\\
&\quad
-
k\int_{\Omega}u^{2}\langle\nabla V,\nabla\varphi\rangle dv.
\end{align}

Substituting \eqref{eqstar} into \eqref{eq46}, we obtain
\begin{align}\label{eqtwostar}
\int_{\partial\Omega}VP\varphi_{\nu}d\sigma
&=
-\int_{\Omega}VPdv
-nk\int_{\Omega}VP\varphi dv
+
\int_{\Omega}P\langle\nabla V,\nabla\varphi\rangle dv
+
\frac{2}{n}\int_{\Omega}V\langle\nabla u,\nabla\varphi\rangle dv
\nonumber\\
&\quad
+
2k\int_{\Omega}uV\langle\nabla u,\nabla\varphi\rangle dv
+
2\int_{\partial\Omega}Vu\nabla^{2}u(\nabla\varphi,\nu) d\sigma
\nonumber\\
&\quad
-
2\int_{\Omega}u\nabla^{2}u(\nabla V,\nabla\varphi)dv
-
2\int_{\Omega}Vu(\operatorname{Ric}-nkg)(\nabla u,\nabla\varphi)dv
\nonumber\\
&\quad
+
\frac{2}{n}\int_{\partial\Omega}Vu\,\varphi_{\nu}d\sigma
-
\frac{2}{n}\int_{\Omega}u\langle\nabla V,\nabla\varphi\rangle dv
-
\frac{2}{n}\int_{\Omega}V\langle\nabla u,\nabla\varphi\rangle dv
\nonumber\\
&\quad
+
2k\int_{\partial\Omega}Vu^{2}\varphi_{\nu}d\sigma
-
4k\int_{\Omega}Vu\langle\nabla u,\nabla\varphi\rangle dv
\nonumber\\
&\quad
-
2k\int_{\Omega}u^{2}\langle\nabla V,\nabla\varphi\rangle dv.
\end{align}

After collecting similar terms, this becomes
\begin{align}
\int_{\partial\Omega}VP\varphi_{\nu} d\sigma
&=
-\int_{\Omega}VP dv
+
\int_{\Omega}P\langle\nabla V,\nabla\varphi\rangle dv
+
2\int_{\partial\Omega}Vu\nabla^{2}u(\nabla\varphi,\nu) d\sigma
\nonumber\\
&\quad
-
2k\int_{\Omega}u^{2}\langle\nabla V,\nabla\varphi\rangle dv
-
2\int_{\Omega}Vu(\operatorname{Ric}-(n-1)kg)(\nabla u,\nabla\varphi) dv
\nonumber\\
&\quad
-
2\int_{\Omega}u\nabla^{2}u(\nabla V,\nabla\varphi) dv
-
\frac{2}{n}\int_{\Omega}u\langle\nabla V,\nabla\varphi\rangle dv
\nonumber\\
&\quad
+
\frac{2}{n}\int_{\partial\Omega}Vu\,\varphi_{\nu} d\sigma
+
2k\int_{\partial\Omega}Vu^{2}\varphi_{\nu} d\sigma
-
nk\int_{\Omega}VP\varphi dv.
\label{eqtwostar_simplified}
\end{align}

Substituting \eqref{eqtwostar_simplified} into \eqref{eq44}, we obtain
\begin{align}\label{eqstar44}
\int_{\partial\Omega}P(Vu)_{\nu}d\sigma
&=
\int_{\partial\Omega}PV(u_{\nu}-\varphi_{\nu})d\sigma
+
\int_{\Omega}P\langle\nabla u,\nabla V\rangle dv
+
\int_{\Omega}u\langle\nabla P,\nabla V\rangle dv
\nonumber\\
&\quad
-
\frac{1}{n-1}
\int_{\Omega}uP\left(\frac{1}{2}X(R)+VR\right)dv
-
\int_{\Omega}VP dv
\nonumber\\
&\quad
+
\int_{\Omega}P\langle\nabla V,\nabla\varphi\rangle dv
+
2\int_{\partial\Omega}Vu\nabla^{2}u(\nabla\varphi,\nu) d\sigma
\nonumber\\
&\quad
-
2k\int_{\Omega}u^{2}\langle\nabla V,\nabla\varphi\rangle dv
-
2\int_{\Omega}Vu(\operatorname{Ric}-(n-1)kg)(\nabla u,\nabla\varphi) dv
\nonumber\\
&\quad
-
2\int_{\Omega}u\nabla^{2}u(\nabla V,\nabla\varphi)dv
-
\frac{2}{n}\int_{\Omega}u\langle\nabla V,\nabla\varphi\rangle dv
\nonumber\\
&\quad
+
\frac{2}{n}\int_{\partial\Omega}Vu\,\varphi_{\nu}d\sigma
+
2k\int_{\partial\Omega}Vu^{2}\varphi_{\nu}d\sigma
-
nk\int_{\Omega}VP\varphi dv.
\end{align}

Returning now to \eqref{eq43}, we obtain
\begin{align}\label{eq49}
-\int_{\Omega}Vu\,\Delta P dv
&=
-\int_{\partial\Omega}Vu\,P_{\nu}d\sigma
+
\int_{\partial\Omega}VP(u_{\nu}-\varphi_{\nu})d\sigma
-
\int_{\Omega}P\langle\nabla V,\nabla(u-\varphi)\rangle dv
\nonumber\\
&\quad
+
\int_{\Omega}u\langle\nabla V,\nabla P\rangle dv
+
2\int_{\partial\Omega}Vu\nabla^{2}u(\nabla\varphi,\nu) d\sigma
\nonumber\\
&\quad
-
2\int_{\Omega}u\nabla^{2}u(\nabla V,\nabla\varphi) dv
-
\frac{2}{n}\int_{\Omega}u\langle\nabla V,\nabla\varphi\rangle dv
\nonumber\\
&\quad
+
\frac{2}{n}\int_{\partial\Omega}Vu\,\varphi_{\nu} d\sigma
+
2k\int_{\partial\Omega}Vu^{2}\varphi_{\nu} d\sigma
\nonumber\\
&\quad
-
2k\int_{\Omega}u^{2}\langle\nabla V,\nabla\varphi\rangle dv
+
nk\int_{\Omega}VP(u-\varphi) dv
\nonumber\\
&\quad
-
2\int_{\Omega}Vu(\operatorname{Ric}-(n-1)kg)(\nabla u,\nabla\varphi) dv.
\end{align}

Next, we collect the terms in \eqref{eq49} involving integrals over \(\Omega\) that do not contain \(k\). Denoting their sum by
\[
C
=
-\int_{\Omega}P\langle\nabla V,\nabla(u-\varphi)\rangle dv
+
\int_{\Omega}u\langle\nabla V,\nabla P\rangle dv
-
2\int_{\Omega}u\nabla^{2}u(\nabla V,\nabla\varphi)dv
-
\frac{2}{n}\int_{\Omega}u\langle\nabla V,\nabla\varphi\rangle dv,
\]
and expanding the expression of \(P\), see \eqref{P_func1} and \eqref{P_funct2}, we find
\begin{align}\label{eq410}
C
&=
-\int_{\Omega}|\nabla u|^{2}
\langle\nabla V,\nabla(u-\varphi)\rangle dv
+
2\int_{\Omega}u\nabla^{2}u\bigl(\nabla V,\nabla(u-\varphi)\bigr) dv
\nonumber\\
&\quad
+
k\int_{\Omega}u^{2}\langle\nabla V,\nabla u\rangle dv
+
k\int_{\Omega}u^{2}\langle\nabla V,\nabla\varphi\rangle dv.
\end{align}

We now compute the first term appearing in \eqref{eq410}. To begin with, observe that
\[
\operatorname{div}\bigl(-|\nabla u|^{2}(u-\varphi)\nabla V\bigr)
=
-(u-\varphi)\langle\nabla |\nabla u|^{2},\nabla V\rangle
-
|\nabla u|^{2}\langle\nabla(u-\varphi),\nabla V\rangle
-
|\nabla u|^{2}(u-\varphi)\Delta V.
\]

Integrating over \(\Omega\) and applying the divergence theorem, we obtain
\begin{align}\label{eqaux0}
-\int_{\Omega}|\nabla u|^{2}\langle\nabla V,\nabla(u-\varphi)\rangle dv
&=
-\int_{\partial\Omega}|\nabla u|^{2}(u-\varphi)V_{\nu} d\sigma
+
\int_{\Omega}(u-\varphi)\langle\nabla |\nabla u|^{2},\nabla V\rangle dv
\nonumber\\
&\quad
+
\int_{\Omega}|\nabla u|^{2}(u-\varphi)\Delta V dv.
\end{align}
Using the identity \eqref{Lapl_V}, we conclude that
\begin{align}\label{eqaux0_final}
-\int_{\Omega}|\nabla u|^{2}\langle\nabla V,\nabla(u-\varphi)\rangle dv
&=
-\int_{\partial\Omega}|\nabla u|^{2}(u-\varphi)V_{\nu} d\sigma
+
2\int_{\Omega}(u-\varphi)\nabla^2 u(\nabla u,\nabla V) dv
\nonumber\\
&\quad
-
\frac{1}{2(n-1)}
\int_{\Omega}X(R)|\nabla u|^{2}(u-\varphi)dv
\nonumber\\
&\quad
-
\frac{1}{n-1}
\int_{\Omega}|\nabla u|^{2}(u-\varphi)VR dv.
\end{align}

On the other hand, using once again the Ricci identity \eqref{Ricci_id}, we obtain
\begin{align*}
\operatorname{div}\bigl(u(u-\varphi)\nabla^{2}u(\nabla V)\bigr)
&=
(u-\varphi)\nabla^{2}u(\nabla V,\nabla u)
+
u\nabla^{2}u(\nabla V,\nabla(u-\varphi))
\nonumber\\
&\quad
+
u(u-\varphi)(\operatorname{Ric}-nkg)(\nabla V,\nabla u)
\nonumber\\
&\quad
+
u(u-\varphi)\langle\nabla^{2}u,\nabla^{2}V\rangle.
\end{align*}

Integrating over \(\Omega\) and applying the divergence theorem, we deduce
\begin{align}\label{auxeq1}
\int_{\Omega}u\nabla^{2}u(\nabla V,\nabla(u-\varphi)) dv
&=
\int_{\partial\Omega}u(u-\varphi)\nabla^{2}u(\nabla V,\nu)d\sigma
-
\int_{\Omega}(u-\varphi)\nabla^{2}u(\nabla V,\nabla u)dv
\nonumber\\
&\quad
-
\int_{\Omega}u(u-\varphi)(\operatorname{Ric}-nkg)(\nabla V,\nabla u)dv
-
\int_{\Omega}u(u-\varphi)\langle\nabla^{2}u,\nabla^{2}V\rangle dv.
\end{align}

Since
\[
\nabla^{2}V
=
\mathring{\nabla^{2}V}
+
\frac{\Delta V}{n}\,g,
\]
using \eqref{Lapl_V} we infer that
\begin{align}\label{eqaux1}
\int_{\Omega}u\nabla^{2}u(\nabla V,\nabla(u-\varphi))dv
&=
\int_{\partial\Omega}u(u-\varphi)\nabla^{2}u(\nabla V,\nu)d\sigma
-
\int_{\Omega}(u-\varphi)\nabla^{2}u(\nabla V,\nabla u) dv
\nonumber\\
&\quad
-
\int_{\Omega}u(u-\varphi)(\operatorname{Ric}-nkg)(\nabla V,\nabla u) dv
-
\int_{\Omega}u(u-\varphi)
\langle\nabla^{2}u,\mathring{\nabla^{2}V}\rangle dv
\nonumber\\
&\quad
-
\frac{1}{2n(n-1)}
\int_{\Omega}u(u-\varphi)X(R) dv
-
\frac{1}{n(n-1)}
\int_{\Omega}u(u-\varphi)VR dv
\nonumber\\
&\quad
-
\frac{k}{2(n-1)}
\int_{\Omega}u^{2}(u-\varphi)X(R) dv
-
\frac{k}{n-1}
\int_{\Omega}u^{2}(u-\varphi)VR dv.
\end{align}

Combining \eqref{eq410}, \eqref{eqaux0_final} and \eqref{eqaux1}, we obtain
\begin{align*}
C
&=
-\int_{\partial\Omega}|\nabla u|^{2}(u-\varphi)V_{\nu}d\sigma
+2\int_{\partial\Omega}u(u-\varphi)\nabla^{2}u(\nabla V,\nu)d\sigma
\nonumber\\
&\quad
-2\int_{\Omega}u(u-\varphi)(\operatorname{Ric}-nkg)(\nabla u,\nabla V)dv
-2\int_{\Omega}u(u-\varphi)\langle\nabla^{2}u,\mathring{\nabla^{2}V}\rangle dv
\nonumber\\
&\quad
-\frac{1}{n-1}\int_{\Omega}(u-\varphi)
\left(
\frac{2}{n}uVR
+2kVu^{2}R
+|\nabla u|^{2}VR
\right)dv
\nonumber\\
&\quad
-\frac{1}{n-1}\int_{\Omega}(u-\varphi)
\left(
\frac12 X(R)|\nabla u|^{2}
+\frac{1}{n}uX(R)
+ku^{2}X(R)
\right)dv
+k\int_{\Omega}u^{2}\langle\nabla V,\nabla(u+\varphi)\rangle dv.
\end{align*}

Substituting this expression for \(C\) into \eqref{eq49}, we deduce
\begin{align}\label{eq411}
-\int_{\Omega}Vu\,\Delta P dv
&=
-\int_{\partial\Omega}|\nabla u|^{2}(u-\varphi)V_{\nu} d\sigma
+2\int_{\partial\Omega}u(u-\varphi)\nabla^{2}u(\nabla V,\nu)d\sigma
\nonumber\\
&\quad
-\int_{\partial\Omega}Vu\,P_{\nu}d\sigma
+\int_{\partial\Omega}VP(u_{\nu}-\varphi_{\nu})d\sigma
+2\int_{\partial\Omega}Vu\nabla^{2}u(\nabla\varphi,\nu)d\sigma
\nonumber\\
&\quad
-\frac{2}{n(n-1)}
\int_{\Omega}u(u-\varphi)VR dv
-\frac{2k}{n-1}
\int_{\Omega}Vu^{2}(u-\varphi)R dv
\nonumber\\
&\quad
-2\int_{\Omega}u(u-\varphi)(\operatorname{Ric}-nkg)(\nabla u,\nabla V)dv
+k\int_{\Omega}u^{2}\langle\nabla V,\nabla(u-\varphi)\rangle dv
\nonumber\\
&\quad
-\frac{1}{n-1}
\int_{\Omega}|\nabla u|^{2}(u-\varphi)VR dv
-\frac{1}{2(n-1)}
\int_{\Omega}X(R)|\nabla u|^{2}(u-\varphi)dv
\nonumber\\
&\quad
-\frac{1}{n(n-1)}
\int_{\Omega}u(u-\varphi)X(R)dv
-\frac{k}{n-1}
\int_{\Omega}u^{2}(u-\varphi)X(R) dv
\nonumber\\
&\quad
-2\int_{\Omega}u(u-\varphi)
\langle\nabla^{2}u,\mathring{\nabla^{2}V}\rangle dv
+\frac{2}{n}\int_{\partial\Omega}Vu\,\varphi_{\nu}d\sigma
+2k\int_{\partial\Omega}Vu^{2}\varphi_{\nu}d\sigma
\nonumber\\
&\quad
+nk\int_{\Omega}VP(u-\varphi)dv
-2\int_{\Omega}Vu(\operatorname{Ric}-(n-1)kg)(\nabla u,\nabla\varphi)dv.
\end{align}

We now expand \(P\) and \(P_{\nu}\) in the previous identity to obtain
\begin{align}\label{eq412}
-\int_{\Omega}Vu\,\Delta P dv
&=
-\int_{\partial\Omega}|\nabla u|^{2}(u-\varphi)V_{\nu} d\sigma
+2\int_{\partial\Omega}u(u-\varphi)\nabla^{2}u(\nabla V,\nu) d\sigma
\nonumber\\
&\quad
+\int_{\partial\Omega}V|\nabla u|^{2}(u_{\nu}-\varphi_{\nu}) d\sigma
-2\int_{\partial\Omega}Vu\nabla^{2}u(\nabla(u-\varphi),\nu) d\sigma
\nonumber\\
&\quad
-k\int_{\partial\Omega}Vu^{2}(u-\varphi)_{\nu} d\sigma
-\frac{2}{n(n-1)}
\int_{\Omega}Vu(u-\varphi)(R-n(n-1)k)dv
\nonumber\\
&\quad
+k\int_{\Omega}u^{2}\langle\nabla V,\nabla(u-\varphi)\rangle dv
-2\int_{\Omega}u(u-\varphi)(\operatorname{Ric}-(n-1)kg)(\nabla u,\nabla V)dv
\nonumber\\
&\quad
-2\int_{\Omega}Vu(\operatorname{Ric}-(n-1)kg)(\nabla u,\nabla\varphi) dv
+2k\int_{\Omega}u(u-\varphi)\langle\nabla u,\nabla V\rangle dv
\nonumber\\
&\quad
-\frac{1}{n-1}
\int_{\Omega}V|\nabla u|^{2}(u-\varphi)(R-(n-1)nk) dv
-\frac{2k}{n-1}
\int_{\Omega}Vu^{2}(u-\varphi)
\left(R-\frac{(n-1)nk}{2}\right) dv
\nonumber\\
&\quad
-\frac{1}{2(n-1)}
\int_{\Omega}X(R)|\nabla u|^{2}(u-\varphi)dv
-\frac{1}{n(n-1)}
\int_{\Omega}u(u-\varphi)X(R)dv
\nonumber\\
&\quad
-\frac{k}{n-1}
\int_{\Omega}u^{2}(u-\varphi)X(R)dv
-2\int_{\Omega}u(u-\varphi)
\langle\nabla^{2}u,\mathring{\nabla^{2}V}\rangle dv.
\end{align}

Next, observe that
\[
\operatorname{div}\bigl((u-\varphi)u^{2}\nabla V\bigr)
=
u^{2}\langle\nabla(u-\varphi),\nabla V\rangle
+
(u-\varphi)\langle\nabla u^{2},\nabla V\rangle
+
(u-\varphi)u^{2}\Delta V.
\]
Integrating by parts, we obtain
\begin{align}\label{eqaux4}
k\int_{\Omega}u(u-\varphi)\langle\nabla V,\nabla u\rangle dv
&=
\frac{k}{2}
\int_{\Omega}(u-\varphi)\langle\nabla V,\nabla u^{2}\rangle dv
\nonumber\\
&=
\frac{k}{2}
\Bigg[
\int_{\partial\Omega}(u-\varphi)u^{2}V_{\nu}d\sigma
-
\int_{\Omega}u^{2}\langle\nabla(u-\varphi),\nabla V\rangle dv
-\int_{\Omega}(u-\varphi)u^{2}\Delta V dv
\Bigg]
\nonumber\\
&=
\frac{k}{2}
\Bigg[
\int_{\partial\Omega}(u-\varphi)u^{2}V_{\nu} d\sigma
+
\frac{1}{2(n-1)}
\int_{\Omega}(u-\varphi)u^{2}X(R) dv
\nonumber\\
&\qquad\qquad
+
\frac{1}{n-1}
\int_{\Omega}(u-\varphi)u^{2}VR dv
-
\int_{\Omega}u^{2}\langle\nabla(u-\varphi),\nabla V\rangle dv
\Bigg],
\end{align}
where we used \eqref{Lapl_V}. Therefore,
\begin{align}\label{eq413}
-\int_{\Omega}Vu\,\Delta P dv
&=
\int_{\partial\Omega}|\nabla u|^{2}
\bigl[V(u-\varphi)_{\nu}-(u-\varphi)V_{\nu}\bigr] d\sigma
+
2\int_{\partial\Omega}
u\nabla^{2}u\bigl((u-\varphi)\nabla V-V\nabla(u-\varphi),\nu\bigr)d\sigma
\nonumber\\
&\quad
+
k\int_{\partial\Omega}
u^{2}\bigl[(u-\varphi)V_{\nu}-V(u-\varphi)_{\nu}\bigr]d\sigma
-\frac{2}{n(n-1)}
\int_{\Omega}Vu(u-\varphi)(R-n(n-1)k) dv
\nonumber\\
&\quad
-2\int_{\Omega}u(u-\varphi)
(\operatorname{Ric}-(n-1)kg)(\nabla u,\nabla V)dv
-2\int_{\Omega}Vu
(\operatorname{Ric}-(n-1)kg)(\nabla u,\nabla\varphi)dv
\nonumber\\
&\quad
-\frac{k}{n-1}
\int_{\Omega}Vu^{2}(u-\varphi)(R-(n-1)nk) dv
-\frac{1}{n-1}
\int_{\Omega}V|\nabla u|^{2}(u-\varphi)(R-(n-1)nk)dv
\nonumber\\
&\quad
-\frac{k}{2(n-1)}
\int_{\Omega}u^{2}(u-\varphi)X(R)dv
-\frac{1}{2(n-1)}
\int_{\Omega}X(R)|\nabla u|^{2}(u-\varphi)dv
\nonumber\\
&\quad
-\frac{1}{n(n-1)}
\int_{\Omega}u(u-\varphi)X(R)dv
-2\int_{\Omega}u(u-\varphi)
\langle\nabla^{2}u,\mathring{\nabla^{2}V}\rangle dv.
\end{align}

This completes the proof of \eqref{eqlema}.
\end{proof}

\end{document}
